\documentclass[11pt,a4paper]{article}

\usepackage[utf8]{inputenc}
\usepackage[T1]{fontenc}
\usepackage[english]{babel}
\usepackage{geometry}
\usepackage{booktabs}
\usepackage{enumitem}
\usepackage{graphicx}
\usepackage{xcolor}
\usepackage{hyperref}
\usepackage{fancyhdr}
\usepackage{tcolorbox}
\usepackage{siunitx}
\usepackage{caption}
\usepackage{float}
\usepackage{amsmath}
\usepackage{pgfplots}
\usepackage{pgfplotstable}
\usepackage{tikz}
\usepackage{multirow}

\pgfplotsset{compat=1.18}
\geometry{margin=2.5cm}

\definecolor{stepcolor}{RGB}{41, 128, 185}
\definecolor{warncolor}{RGB}{231, 76, 60}
\definecolor{tipcolor}{RGB}{39, 174, 96}
\definecolor{paramcolor}{RGB}{142, 68, 173}
\definecolor{baseline}{RGB}{52, 73, 94}
\definecolor{bilayer}{RGB}{39, 174, 96}

\newtcolorbox{warningbox}[1][]{colback=warncolor!5,colframe=warncolor!80,fonttitle=\bfseries,title={#1}}
\newtcolorbox{tipbox}[1][]{colback=tipcolor!5,colframe=tipcolor!80,fonttitle=\bfseries,title={#1}}

\pagestyle{fancy}
\fancyhf{}
\fancyhead[L]{\small Yield Analysis Methodology}
\fancyhead[R]{\small EGOFET Fabrication}
\fancyfoot[C]{\thepage}

\title{
    \vspace{-1cm}
    \textbf{Yield Analysis for EGOFET Electrode Patterning} \\
    \large Measuring the Impact of the PMGI/S1813 Bilayer Lift-Off Process
}
\author{
    e-MolMat Group -- ICMAB-CSIC \\
    \texttt{v1.0 -- \today}
}
\date{}

\begin{document}

\maketitle
\thispagestyle{fancy}

\begin{abstract}
This document defines a methodology for quantifying the success rate of
the EGOFET electrode patterning process. It introduces three complementary
metrics: \textbf{visual yield} (microscope inspection), \textbf{electrical
yield} (device functionality), and \textbf{pattern integrity} (structural
quality). These metrics enable a rigorous before/after comparison between
the current single-layer S1813 protocol and the proposed PMGI/S1813
bilayer protocol. A data recording template is provided for systematic
tracking across fabrication batches.
\end{abstract}

\tableofcontents
\newpage

% ============================================================
\section{Motivation}
% ============================================================

The primary failure mode in current EGOFET fabrication is residual gold
in the Source--Drain channels after lift-off. This forces manual removal
with cotton swabs, which damages electrode patterns and reduces
device-to-device reproducibility.

To evaluate whether the PMGI/S1813 bilayer process improves yield, we
need \textbf{quantitative metrics} that capture:
\begin{enumerate}
    \item How many devices survive fabrication without defects.
    \item How many devices are electrically functional.
    \item How uniform the electrode patterns are across a batch.
\end{enumerate}

These metrics also serve as a diagnostic tool: if yield drops in future
batches, the data can pinpoint which fabrication step degraded.

% ============================================================
\section{Definitions}
% ============================================================

\subsection{Device}

A \textbf{device} is a single EGOFET structure on the Kapton substrate,
consisting of:
\begin{itemize}
    \item Source electrode (Au)
    \item Drain electrode (Au)
    \item Channel region (between Source and Drain)
    \item Gate electrode (Au, coplanar)
    \item Gate area (defined by dextran or geometry)
\end{itemize}

Each substrate typically contains multiple devices (e.g., 4--8 per chip).
All devices on a substrate are counted individually.

\subsection{Batch}

A \textbf{batch} is a single fabrication run, identified by date and
substrate number. Example: \texttt{2026-09-15\_S03} = September 15, 2026,
Substrate 3.

% ============================================================
\section{Yield Metrics}
% ============================================================

\subsection{Metric 1: Visual Yield ($Y_v$)}

\textbf{Definition:} Percentage of devices where the channel region is
\textbf{completely free} of gold residue after lift-off, as determined
by optical microscope inspection.

\begin{equation}
    Y_v = \frac{N_{\text{clean channels}}}{N_{\text{total devices}}}
    \times 100\%
\end{equation}

\begin{tipbox}[Inspection Protocol]
\begin{enumerate}
    \item Inspect each device at $\geq 20\times$ magnification after
          lift-off and cleaning.
    \item A device is \textbf{PASS} if: no gold islands, no metal fences,
          no visible debris in the channel region.
    \item A device is \textbf{FAIL} if: any gold residue is visible in
          the channel, regardless of size.
    \item Record the number of PASS and FAIL devices per substrate.
\end{enumerate}
\end{tipbox}

\subsection{Metric 2: Pattern Integrity ($Y_p$)}

\textbf{Definition:} Percentage of devices where the gold electrode
pattern is \textbf{structurally intact} after lift-off --- i.e., no
delamination, no scratching, no broken electrodes.

\begin{equation}
    Y_p = \frac{N_{\text{intact patterns}}}{N_{\text{total devices}}}
    \times 100\%
\end{equation}

\begin{tipbox}[Inspection Protocol]
\begin{enumerate}
    \item Inspect each device at $\geq 10\times$ magnification.
    \item A device is \textbf{PASS} if: all electrodes are continuous,
          edges are sharp, no peeling or cracking visible.
    \item A device is \textbf{FAIL} if: any electrode is broken, peeled,
          scratched, or has irregular edges.
    \item Record the number of PASS and FAIL devices per substrate.
\end{enumerate}
\end{tipbox}

\subsection{Metric 3: Electrical Yield ($Y_e$)}

\textbf{Definition:} Percentage of devices that show \textbf{functional}
transistor behavior after semiconductor deposition and gating.

\begin{equation}
    Y_e = \frac{N_{\text{functional devices}}}{N_{\text{total devices}}}
    \times 100\%
\end{equation}

\begin{tipbox}[Functional Criteria]
A device is \textbf{functional} if all of the following are met:
\begin{enumerate}
    \item Transfer curve ($I_D$ vs.\ $V_G$) shows clear ON/OFF switching.
    \item $I_{\text{ON}} / I_{\text{OFF}} \geq 10$.
    \item Threshold voltage $|V_{\text{th}}| < \SI{1}{\V}$.
    \item Hysteresis between forward and reverse sweeps
          $\Delta V_{\text{th}} < \SI{0.3}{\V}$.
\end{enumerate}
\end{tipbox}

\subsection{Composite Yield Score}

For overall process comparison, define a \textbf{composite yield}:

\begin{equation}
    Y_{\text{composite}} = Y_v \times Y_p \times Y_e
\end{equation}

This penalizes any failure mode. For example, if $Y_v = 95\%$,
$Y_p = 98\%$, and $Y_e = 90\%$, then $Y_{\text{composite}} = 83.7\%$.

% ============================================================
\section{Experimental Design}
% ============================================================

\subsection{Before/After Comparison}

To demonstrate the effectiveness of the bilayer process, perform a
controlled comparison:

\begin{table}[H]
\centering
\caption{Experimental design for before/after comparison.}
\begin{tabular}{@{}lcc@{}}
\toprule
\textbf{Parameter} & \textbf{Group A (Baseline)} & \textbf{Group B (Bilayer)} \\
\midrule
Photoresist & S1813 only & PMGI SF5 + S1813 \\
Lift-off solvent & Acetone + US & DMSO at \SI{60}{\celsius} \\
Number of substrates & $\geq 3$ & $\geq 3$ \\
Devices per substrate & $\geq 4$ & $\geq 4$ \\
Total devices per group & $\geq 12$ & $\geq 12$ \\
Evaporation batch & Same & Same \\
OSC deposition & Same & Same \\
Measurement protocol & Same & Same \\
\bottomrule
\end{tabular}
\end{table}

\begin{warningbox}[Control Variables]
To ensure a fair comparison, \textbf{all other parameters must be
identical} between groups A and B:
\begin{itemize}
    \item Same Kapton batch and cutting procedure.
    \item Same Si carrier wafers (or equivalent).
    \item Same evaporator run (if possible) or same technician with same
          parameters.
    \item Same OSC ink batch and BAMS parameters.
    \item Same SAM protocol (PFBT concentration, immersion time).
    \item Same measurement equipment and settings.
\end{itemize}
\end{warningbox}

\subsection{Sample Size Justification}

With $n = 12$ devices per group and a baseline yield of $Y_0 = 50\%$
(estimated from current process), we can detect a $\geq 25$ percentage
point improvement ($Y_1 = 75\%$) with:
\begin{itemize}
    \item Significance level $\alpha = 0.05$
    \item Power $1 - \beta = 0.80$
\end{itemize}

This is a conservative estimate. If baseline yield is lower (e.g., 30\%),
even smaller improvements become statistically significant.

% ============================================================
\section{Data Recording}
% ============================================================

\subsection{Per-Device Record}

For each device, record the following fields in the data template
(see Appendix~\ref{app:template}):

\begin{table}[H]
\centering
\caption{Per-device data fields.}
\begin{tabular}{@{}llp{6cm}@{}}
\toprule
\textbf{Field} & \textbf{Type} & \textbf{Description} \\
\midrule
\texttt{batch\_id} & String & Batch identifier (e.g., \texttt{2026-09-15\_S03}) \\
\texttt{device\_id} & Integer & Device number within substrate \\
\texttt{process} & String & \texttt{baseline} or \texttt{bilayer} \\
\texttt{date} & Date & Fabrication date \\
\texttt{visual\_pass} & Boolean & Channel clean after lift-off? \\
\texttt{pattern\_pass} & Boolean & Electrode pattern intact? \\
\texttt{cotton\_swab} & Boolean & Cotton swab used for cleanup? \\
\texttt{notes} & String & Free-text observations \\
\texttt{Vth} & Float & Threshold voltage (V) \\
\texttt{Ion\_Ioff} & Float & ON/OFF current ratio \\
\texttt{mobility} & Float & Field-effect mobility (cm$^2$/Vs) \\
\texttt{hysteresis} & Float & $\Delta V_{\text{th}}$ (V) \\
\texttt{functional} & Boolean & Meets functional criteria? \\
\bottomrule
\end{tabular}
\end{table}

\subsection{Batch Summary}

After each batch, compute:

\begin{table}[H]
\centering
\caption{Batch summary fields.}
\begin{tabular}{@{}ll@{}}
\toprule
\textbf{Metric} & \textbf{Formula} \\
\midrule
$Y_v$ & \texttt{visual\_pass} count / total devices \\
$Y_p$ & \texttt{pattern\_pass} count / total devices \\
$Y_e$ & \texttt{functional} count / total devices \\
$Y_{\text{composite}}$ & $Y_v \times Y_p \times Y_e$ \\
Cotton swab rate & \texttt{cotton\_swab} count / total devices \\
\bottomrule
\end{tabular}
\end{table}

% ============================================================
\section{Expected Outcomes}
% ============================================================

Based on literature and the physics of the bilayer undercut:

\begin{table}[H]
\centering
\caption{Expected yield improvement.}
\begin{tabular}{@{}lcc@{}}
\toprule
\textbf{Metric} & \textbf{Baseline (S1813)} & \textbf{Bilayer (PMGI+S1813)} \\
\midrule
$Y_v$ (visual) & 40--60\% & \textbf{90--100\%} \\
$Y_p$ (pattern) & 60--80\% & \textbf{95--100\%} \\
$Y_e$ (electrical) & 30--50\% & \textbf{80--95\%} \\
Cotton swab rate & 40--60\% & \textbf{$<$5\%} \\
\bottomrule
\end{tabular}
\end{table}

\begin{tipbox}[Interpreting Results]
\begin{itemize}
    \item If $Y_v$ improves but $Y_e$ does not: the problem may be in
          OSC deposition or SAM, not in lift-off.
    \item If $Y_p$ improves but $Y_v$ does not: the undercut may be
          insufficient --- increase develop time.
    \item If cotton swab rate drops to near zero: the bilayer is working
          as intended.
\end{itemize}
\end{tipbox}

% ============================================================
\section{Visualization}
% ============================================================

The data should be visualized as:

\begin{enumerate}
    \item \textbf{Bar chart}: $Y_v$, $Y_p$, $Y_e$ side by side for
          baseline vs.\ bilayer groups.
    \item \textbf{Time series}: Yield metrics across batches to track
          process stability.
    \item \textbf{Box plots}: Distribution of $V_{\text{th}}$, mobility,
          and $I_{\text{ON}}/I_{\text{OFF}}$ for each group.
\end{enumerate}

Example visualization code is provided in the companion Jupyter notebook
(\texttt{notebooks/yield\_analysis.ipynb}).

% ============================================================
\section{References}
% ============================================================

\begin{enumerate}[label={[\arabic*]}]
    \item Ruiz-Molina, S. PhD Thesis, ICMAB-CSIC, 2024.
    \item Ricci, S. PhD Thesis, ICMAB-CSIC, 2020.
    \item MicroChem Corp. \textit{PMGI Resists --- Technical Data Sheet}.
    \item Fraunhofer IPMS. ``OFET Substrates --- Test Chips for Material
          Characterization.'' Protocol, 2026.
\end{enumerate}

% ============================================================
\appendix
\section{Data Recording Template}
\label{app:template}
% ============================================================

The complete data template is maintained as a structured file in the
\texttt{data/} directory. See \texttt{data/yield\_template.md} for the
field definitions and \texttt{data/raw/} for individual batch records.

\subsection{Template Format}

Each batch record is a YAML file named
\texttt{data/raw/\{batch\_id\}.yaml} with the following structure:

\begin{verbatim}
batch_id: "2026-09-15_S03"
date: "2026-09-15"
process: "bilayer"          # "baseline" or "bilayer"
substrate: "Kapton 75um"
photoresist: "PMGI SF5 + S1813"
lift_off: "DMSO 60C, 2.5h"
evaporation:
  Cr_thickness_nm: 5
  Au_thickness_nm: 40
  rate_Cr: 0.2              # A/s
  rate_Au: 0.5              # A/s
devices:
  - device_id: 1
    visual_pass: true
    pattern_pass: true
    cotton_swab: false
    Vth: -0.15
    Ion_Ioff: 45
    mobility: 0.8
    hysteresis: 0.12
    functional: true
    notes: ""
  - device_id: 2
    visual_pass: true
    pattern_pass: true
    cotton_swab: false
    Vth: -0.22
    Ion_Ioff: 38
    mobility: 0.7
    hysteresis: 0.18
    functional: true
    notes: "minor edge roughness"
\end{verbatim}

\subsection{Batch Summary Computation}

After recording all devices, compute the batch summary:

\begin{verbatim}
batch_summary:
  total_devices: 8
  visual_pass_count: 8
  pattern_pass_count: 8
  cotton_swab_count: 0
  functional_count: 7
  Yv: 100.0                # %
  Yp: 100.0                # %
  Ye: 87.5                 # %
  Ycomposite: 87.5         # %
  mean_Vth: -0.18
  std_Vth: 0.04
  mean_mobility: 0.75
  mean_Ion_Ioff: 41
\end{verbatim}

\end{document}
